\chapter{Plan de validacion en CoppeliaSim}

El uso de CoppeliaSim se plantea como una fase de reduccion de riesgo antes de cualquier prueba fisica. La simulacion no certifica la seguridad industrial, pero permite comprobar coherencia de rutas, maniobrabilidad, aproximacion a estaciones, evitacion de obstaculos y logica de misiones con una base omnidireccional tipo KUKA YouBot, KUKA Omnirob o equivalente.

\section{Modelo de simulacion}

\begin{longtable}{P{0.24\textwidth}P{0.62\textwidth}}
\toprule
\textbf{Elemento simulado} & \textbf{Descripcion}\\
\midrule
\endhead
Robot & Base omnidireccional con volumen equivalente a AMR industrial ligero, velocidad limitada y sensores virtuales.\\
Mapa de planta & Pasillos, estaciones de HTP, zona de preparacion de kits, punto de carga y zona de retorno.\\
Obstaculos & Utiles, carros, estructuras, operarios simulados como objetos dinamicos y zonas temporalmente bloqueadas.\\
Sensores & LiDAR 2D/3D para obstaculos, camara RGB-D conceptual para FOD y lectores RFID/QR representados por puntos de validacion.\\
Misiones & Entrega de kit, entrega de herramienta, retorno de util, ronda FOD y mision cancelada por obstaculo persistente.\\
\bottomrule
\end{longtable}

\section{Escenarios de prueba}

\begin{longtable}{P{0.25\textwidth}P{0.39\textwidth}P{0.26\textwidth}}
\toprule
\textbf{Escenario} & \textbf{Prueba} & \textbf{Resultado esperado}\\
\midrule
\endhead
Navegacion base & Ruta desde preparacion de kits hasta estacion HTP. & Llegada dentro de tolerancia sin colision.\\
Aproximacion lateral & Entrada a estacion con maniobra omnidireccional. & Posicionamiento sin giros amplios ni bloqueo de pasillo.\\
Obstaculo dinamico & Operario o carro cruza la trayectoria. & Reduccion de velocidad, parada o replanificacion.\\
Ruta bloqueada & Pasillo cerrado por util temporal. & Recalculo de ruta o mision en espera.\\
Entrega trazada & Lectura RFID/QR en carga y destino. & Registro completo de kit, estacion y hora.\\
Ronda FOD & Captura conceptual de zona antes/despues de entrega. & Evento registrado con imagen o marca de inspeccion.\\
Multi-robot & Dos o mas AMR comparten pasillo. & Sin interbloqueos; prioridad definida por gestor de flota.\\
\bottomrule
\end{longtable}

\section{Criterios de salida de simulacion}

Antes de pasar a piloto fisico, el modelo debe demostrar que las rutas propuestas no invaden zonas restringidas, que el robot puede aproximarse a estaciones sin maniobras excesivas, que las paradas de seguridad no bloquean el flujo de forma recurrente y que las misiones generan registros completos. Los parametros finales de simulacion deben guardarse como anexo tecnico: mapa, velocidad maxima, radios de seguridad, tolerancia de posicion, zonas prohibidas y tiempos de mision.
